% Generated draft: requires mathematical and source review.
\documentclass[11pt]{article}
\usepackage[margin=25mm]{geometry}
\usepackage{fontspec}
\setmainfont{Latin Modern Roman}
\usepackage{amsmath,amssymb,mathtools,unicode-math}
\setmathfont{Latin Modern Math}
\usepackage{enumitem,needspace,xurl,hyperref}
\setlength{\parindent}{0pt}
\setlength{\parskip}{5pt}
\setlength{\emergencystretch}{3em}
\newcommand{\N}{\mathbb{N}}
\newcommand{\Z}{\mathbb{Z}}
\newcommand{\Q}{\mathbb{Q}}
\newcommand{\R}{\mathbb{R}}
\newcommand{\C}{\mathbb{C}}
\newcommand{\F}{\mathbb{F}}
\newcommand{\A}{\mathbb{A}}
\newcommand{\PP}{\mathbb{P}}
\newcommand{\E}{\mathbb{E}}
\newcommand{\eps}{\varepsilon}
\newcommand{\dd}{\,\mathrm{d}}
\begin{document}
\section*{UnsolvedMath ID: 30002714}
\textbf{Source problem code:} catalogue entry
\subsection{Definitions and mathematical statement}
For a commutative integral domain \(A\) with fraction field \(K=\operatorname{Frac}(A)\), its \emph{complete integral closure in \(K\)} is
\[
A^{\mathrm{cic}}=\left\{x\in K:\text{there exists }0\ne a\in A\text{ such that }ax^n\in A\text{ for every integer }n\geq 1\right\}.
\]
A subring \(W\subseteq K\) is a \emph{valuation ring of \(K\)} if, for every \(x\in K\setminus\{0\}\), either \(x\in W\) or \(x^{-1}\in W\).

The supplied question is: for every total blow-up ring
\[
S=B_{\infty}(R,v),
\]
is the complete integral closure of \(S\) a valuation ring? Under the usual integral-domain interpretation, with \(K=\operatorname{Frac}(S)\) and complete integral closure taken in \(K\), the requested assertion is
\[
\forall x\in K\setminus\{0\},\qquad x\in S^{\mathrm{cic}}\ \text{or}\ x^{-1}\in S^{\mathrm{cic}}.
\]
The supplied record does not define \(B_{\infty}(R,v)\), specify the admissible rings \(R\) and valuations \(v\), or state the ambient ring or field for complete integral closure. Consequently, the displayed assertion is a conditional formulation using the usual domain convention; it does not supply the missing definition of a total blow-up ring or establish that this convention matches the source. The intended quantifier ranges over all total blow-up rings allowed by that source.
\subsection{Short English statement}
The question asks whether taking the complete integral closure of any total blow-up ring always produces a valuation ring. The construction of a total blow-up ring and its hypotheses are absent from the supplied record, so the precise class of rings remains unresolved. The supplied open status is dated curation, not an independently verified conclusion.
\subsection{Sources}
\begin{itemize}
\item Supplied local registry record: permanent UnsolvedMath ID \(30002714\), title \emph{Complete Integral Closures of Total Blow-Up Rings}, source problem code \texttt{catalogue entry}, in the Oberwolfach Reports Open Problems collection. The proposed year is recorded as 2014; no individual proposer is supplied.
\item Supplied source citation: \emph{Valuation Theory and Its Applications} (2015), Oberwolfach Reports, pp.\ 2757--2823, DOI \texttt{10.4171/owr/2014/49}. \url{https://doi.org/10.4171/owr/2014/49}. The background describes Oberwolfach Reports as documenting mathematical workshops and their open-problem discussions.
\item Supplied literature source: \emph{Oberwolfach Report 49/2014}. \url{https://ems.press/content/serial-article-files/46540}.
\item Supplied literature source: Heinzer, \emph{Some remarks on complete integral closure}. \url{https://doi.org/10.1017/S1446788700007230}.
\item Literature assessment dated exactly \texttt{2026-08-21}: the supplied curation labels the problem open, states that the question for total blow-up rings is explicitly left open, and reports that no later general resolution was found. This is supplied dated curation derived from the OWR material, not an independently verified literature assessment; time-sensitive status and cited claims require independent verification.
\end{itemize}
\end{document}
